\documentclass[journal]{IEEEtran}

\usepackage{bm}
\usepackage{bbm}
\usepackage{amsmath}
\usepackage{mathtools}
\usepackage{cuted}
\usepackage{algorithmic}
\usepackage{algorithm}
\usepackage{cite}
\usepackage{color}
\usepackage{mathrsfs}
\usepackage{amsfonts}
\usepackage{amssymb}
\usepackage{makecell}
\usepackage{multirow}
\usepackage{times}
\usepackage{makecell}
\usepackage{tabulary}
\usepackage{array}
\usepackage{url}
\usepackage[table]{xcolor}
\usepackage{colortbl}  
\usepackage{float}
\usepackage{balance}
\usepackage{MnSymbol}
\usepackage{setspace}
\usepackage{booktabs}
\usepackage[normalem]{ulem}
\usepackage{tikz}
\usepackage{xcolor}
\usepackage{graphicx}
\usepackage{booktabs}
\usepackage{diagbox}
% \usepackage[flushleft]{threeparttable}
\usepackage{threeparttable}
\usepackage[caption=false,font=normalsize,labelfont=sf,textfont=sf]{subfig}
\usepackage{multirow}
\usepackage{adjustbox}
\usepackage[table]{xcolor}
\usepackage{colortbl}
\usepackage{pifont}
\usepackage{tikz}

\newcommand{\legendbox}[1]{%
  \tikz\draw[line width=0.8pt, draw=#1, fill=none] 
    (0,0) rectangle (0.6em,0.6em);%
}
% 定义颜色，使用 RGB 值，并归一化到 [0,1] 范围内
\definecolor{lightblue}{RGB}{97,203,244}
\definecolor{darkred}{RGB}{192,0,0}
\definecolor{lightgreen}{RGB}{142,217,115}
\newtheorem{theorem}{Theorem}
\newtheorem{definition}{Definition}
\newtheorem{assumption}{Assumption}
\newtheorem{remark}{Remark}


% *** GRAPHICS RELATED PACKAGES ***
%
\ifCLASSINFOpdf
  % \usepackage[pdftex]{graphicx}
  % declare the path(s) where your graphic files are
  % \graphicspath{{../pdf/}{../jpeg/}}
  % and their extensions so you won't have to specify these with
  % every instance of \includegraphics
  % \DeclareGraphicsExtensions{.pdf,.jpeg,.png}
\else
  % or other class option (dvipsone, dvipdf, if not using dvips). graphicx
  % will default to the driver specified in the system graphics.cfg if no
  % driver is specified.
  % \usepackage[dvips]{graphicx}
  % declare the path(s) where your graphic files are
  % \graphicspath{{../eps/}}
  % and their extensions so you won't have to specify these with
  % every instance of \includegraphics
  % \DeclareGraphicsExtensions{.eps}
\fi




% % correct bad hyphenation here
% \hyphenation{op-tical net-works semi-conduc-tor}


\begin{document}
%
% paper title
% Titles are generally capitalized except for words such as a, an, and, as,
% at, but, by, for, in, nor, of, on, or, the, to and up, which are usually
% not capitalized unless they are the first or last word of the title.
% Linebreaks \\ can be used within to get better formatting as desired.
% Do not put math or special symbols in the title.
\title{\huge
SCOPE: Surround-view Collision Prediction via Panoramic Oriented TTC Estimation
}
%
%
% author names and IEEE memberships
% note positions of commas and nonbreaking spaces ( ~ ) LaTeX will not break
% a structure at a ~ so this keeps an author's name from being broken across
% two lines.
% use \thanks{} to gain access to the first footnote area
% a separate \thanks must be used for each paragraph as LaTeX2e's \thanks
% was not built to handle multiple paragraphs
%

% \author{Chunyu Feng, Yeqiang Qian, Ming Yang
% % \thanks{*This work is supported by the National Natural Science Foundation of China under Grant 62103261. (Corresponding author: Yeqiang Qian and Ming Yang.)}% <-this % stops a space
% % \thanks{Nan Ming, Yeqiang Qian, Chunyu Feng, Chunxiang Wang and Ming Yang are with the Department of Automation, Shanghai Jiao Tong University, Key Laboratory of System Control and Information Processing, Ministry of Education of China, Shanghai, 200240, China (email: {\tt\small qianyeqiang@sjtu.edu.cn}).}
% }


% The paper headers
\markboth{IEEE TRANSACTIONS ON ROBOTICS}%
{Shell \MakeLowercase{\textit{et al.}}: Bare Demo of IEEEtran.cls for IEEE Journals}
% The only time the second header will appear is for the odd numbered pages
% after the title page when using the twoside option.
% 
% *** Note that you probably will NOT want to include the author's ***
% *** name in the headers of peer review papers.                   ***
% You can use \ifCLASSOPTIONpeerreview for conditional compilation here if
% you desire.




% If you want to put a publisher's ID mark on the page you can do it like
% this:
%\IEEEpubid{0000--0000/00\$00.00~\copyright~2015 IEEE}
% Remember, if you use this you must call \IEEEpubidadjcol in the second
% column for its text to clear the IEEEpubid mark.



% use for special paper notices
%\IEEEspecialpapernotice{(Invited Paper)}




% make the title area
\maketitle

% As a general rule, do not put math, special symbols or citations
% in the abstract or keywords.
\begin{abstract}
Omnidirectional collision avoidance is a prerequisite for the safe deployment of autonomous mobile robots. While vision-based Time-to-Collision (TTC) methods provide robust collision risk assessment, they are typically limited to monocular forward views and suffer from directional ambiguity.
In this paper, we present SCOPE, a panoramic framework that delivers 360° collision risk assessment from surround-view cameras. To achieve this, we design a unified Range View (RV) representation that projects multi-view features into a seamless, continuous domain. Crucially, we introduce Oriented TTC (OTTC), a vector-based metric that combines collision urgency with relative motion direction, enabling the system to ranks collision risks accurately across all directions.
To address the scarcity of training data, we also reconstruct a densified \textit{nuScenes-SceneFlow} dataset from raw point clouds to enable pixel-level supervision. Our method achieves state-of-the-art performance on public benchmarks and demonstrates robust zero-shot transfer capabilities on a real vehicle platform. 
The code and data will be available upon the publication of this paper.

% Omnidirectional collision warning is fundamental to practical deployment in robotic systems.
% Vision-based time-to-collision (TTC) offers category-agnostic collision cues compared to detection pipelines, enhancing robustness to novel agents and long-tail objects.
% But existing TTC estimation methods are largely monocular and forward-view only, leaving lateral and rear hazards insufficiently covered. 
% To address this limitation, we propose SCOPE, an end-to-end panoramic TTC framework that delivers 360° risk assessment from surround-view cameras.
% SCOPE adopt a range view (RV) unified representation that maps all cameras to a common spherical, ray-consistent domain. This preserves angular continuity across seams, removes duplicate observations in overlaps, and yields a single, consistent TTC value for every viewing direction. 
% To reduce false alarms and correctly ranks hazards across 360°, we further introduce orientation-aware TTC that estimates the direction of relative motion and combines it with TTC to form a joint risk representation.
% Because existing surround-view datasets such as nuScenes lack dense, pixel-level supervision for TTC, we reconstruct a denser \textit{nuScenes-SceneFlow} dataset from the raw nuScenes point clouds and annotations, and derive per-pixel scale and orientation ground truth from it. 
% Extensive experiments on the public dataset demonstrate that our method achieves state-of-the-art performance over monocular TTC baselines and additionally supports zero-shot, real-time operation on a local autonomous-driving platform. 
% Our code is available at \url{https://github.com/ChunyuFeng/SCOPE}.
\end{abstract}

% Note that keywords are not normally used for peerreview papers.
\begin{IEEEkeywords}
Omnidirectional collision avoidance, oriented time-to-collision, range view representation, surround-view cameras, autonomous mobile robots.
\end{IEEEkeywords}
% For peer review papers, you can put extra information on the cover
% page as needed:
% \ifCLASSOPTIONpeerreview
% \begin{center} \bfseries EDICS Category: 3-BBND \end{center}
% \fi
%
% For peerreview papers, this IEEEtran command inserts a page break and
% creates the second title. It will be ignored for other modes.
\IEEEpeerreviewmaketitle

%%%%%%%%%%%%%%%%%%%%%%%%%%%%%%%%%%%%%%%%%%%%%%%%%%%%%%%%%%%%%%%%%%%%%%%%%%%%%%%%

\section{Introduction}
\label{sec:intro}


\IEEEPARstart{O}{bstacle avoidance} is a fundamental capability for autonomous mobile robots, serving as a critical safeguard for downstream tasks such as path planning and navigation~\cite{francis2020long, zhang2025end, lu2022online}. The core of obstacle avoidance lies in collision prediction, i.e., assessing whether and when the robot will collide with surrounding objects. A key metric in this process is Time-to-Collision (TTC), which measures the remaining time before impact and directly encodes the urgency of potential collisions.

Existing approaches to collision prediction largely follow two paradigms. 
Traditional perception pipelines construct explicit scene representations via object detection~\cite{pan2021pointformer,wang2022detr3d,li2024bevformer}, 3D semantic occupancy~\cite{wei2023surroundocc,yang2024adaptiveocc,tang2024sparseocc}, and depth estimation~\cite{schmied2023r3d3,xie2023omnividar,deng2025omnistereo}. 
These intermediate outputs are then post-processed to derive TTC. 
An alternative paradigm eliminates these intermediate steps. Direct TTC estimation methods~\cite{yang2020opticalexpansion,badki2021binaryttc,li2024fpttc} exploit optical expansion, where the rate of image-scale change directly reflects the rate of depth change. 
By leveraging this geometric principle, such systems infer collision risks from visual inputs directly, sidestepping the semantic constraints of detection-based models and the scale ambiguity inherent in monocular depth estimation.


% \IEEEPARstart{O}{bstacle avoidance} is a fundamental capability for autonomous mobile robots. It serves as a critical safeguard for downstream tasks such as path planning and navigation.
% Traditional perception pipelines address this requirement by constructing explicit scene representations, utilizing modules such as object detection~\cite{pan2021pointformer,wang2022detr3d,li2024bevformer}, 3D semantic occupancy~\cite{wei2023surroundocc,yang2024adaptiveocc,tang2024sparseocc}, and depth estimation~\cite{schmied2023r3d3,xie2023omnividar,deng2025omnistereo}. These intermediate outputs are then post-processed to derive Time-to-Collision (TTC), the metric quantifying the urgency of potential impacts.
% An alternative paradigm eliminates these intermediate steps. Direct TTC estimation methods~\cite{yang2020opticalexpansion,badki2021binaryttc,li2024fpttc} exploit optical expansion, where the rate of image scale change directly reflects the rate of depth change. By using this geometric principle, such systems infer collision risks end-to-end from visual inputs. This approach successfully sidesteps the semantic constraints of detection-based models and the scale ambiguity inherent in monocular depth estimation.

However, current direct TTC estimation frameworks face two fundamental limitations: \textbf{spatial incompleteness} and \textbf{directional ambiguity}. Regarding \textbf{spatial incompleteness}, existing paradigms are predominantly monocular and restricted to the forward field of view~\cite{badki2021binaryttc,li2024fpttc}. Forward perception only is inadequate for highly dynamic scenarios, where hazards can emerge unpredictably from any direction surrounding the robot as illustrated in Fig.~\ref{fig:limitations} (a).

Even more critical is the \textbf{directional ambiguity}. TTC reflects only the motion along the line of sight between the camera and the obstacle, ignoring lateral motion components. This ambiguity is particularly problematic in omnidirectional scenarios, as parallel motion without collision risk is extremely common. 
% For instance, vehicles overtaking in adjacent lanes may generate a small TTC value indicating high urgency, even though it poses no actual collision risk. Consequently, relying solely on TTC results will lead to frequent false alarms.
As Fig.~\ref{fig:limitations} (c) illustrates, two vehicles passing side by side may yield a small TTC value which indicating an immediate collision. These false alarms can lead to unnecessary emergency maneuvers, undermining system reliability.

\begin{figure}[!t]  
    \centering
    \includegraphics[width=0.45\textwidth]{figs/RV.pdf}
    \caption{
        % Comparison between the discrete surround-view representation and unified RV representation. The outer hexagon displays six separate images captured by cameras surrounding the ego vehicle. Conversely, the inner colored cylinders depict panoramic RV representation that projects the scene around the vehicle in a continuous domain, providing a unified encoding across all directions.
        Comparison between discrete surround-view inputs and the unified OTTC prediction. The outer hexagon displays six separate camera images. In contrast, the inner cylinders present the proposed OTTC in a continuous Range View. This representation combines scalar TTC and collision orientation into a seamless panoramic field, resolving the overlap issue of multi-view inputs.}
    \label{fig:range_view}
\end{figure}

\begin{figure*}[t]
    \centering
    \begin{minipage}{0.98\textwidth}
        \centering
        \begin{minipage}{0.48\textwidth}
            \centering
            \includegraphics[width=\linewidth]{figs/Spatial_Incompleteness0.pdf}
            % \\[-0.3em]
            \small (a) 
        \end{minipage}
        \hfill
        \begin{minipage}{0.48\textwidth}
            \centering
            \includegraphics[width=\linewidth]{figs/Spatial_Incompleteness1.pdf}
            % \\[-0.3em]
            \small (b) 
        \end{minipage}
    \end{minipage}

    \smallskip

    \begin{minipage}{0.98\textwidth}
        \centering
        \begin{minipage}{0.48\textwidth}
            \centering
            \includegraphics[width=\linewidth]{figs/Direction_Ambiguity0.pdf}
            % \\[-0.3em]
            \small (c) 
        \end{minipage}
        \hfill
        \begin{minipage}{0.48\textwidth}
            \centering
            \includegraphics[width=\linewidth]{figs/Direction_Ambiguity1.pdf}
            % \\[-0.3em]
            \small (d) 
        \end{minipage}
    \end{minipage}

    \caption{Illustration of the two fundamental limitations of traditional forward-view TTC estimation methods and how surround-view OTTC addresses them.
    (a) Spatial incompleteness: forward TTC leaves large blind regions, where vehicles can pose undetected collision risks. 
    (b) A surround-view camera suite removes TTC blind regions and enables omnidirectional collision awareness. 
    (c) Directional ambiguity: TTC is computed only from radial motion, ignoring the lateral component of relative motion, so a parallel pass-by can yield a small TTC and trigger a false alarm. 
    (d) By jointly considering TTC and the orientation of relative motion, it correctly recognizes the scenario as a safe pass-by.}
    \label{fig:limitations}
\end{figure*}

We propose \textbf{SCOPE} (\textbf{S}urround-view \textbf{CO}llision \textbf{P}rediction via oriented ttc \textbf{E}stimation), a panoramic collision prediction framework designed to overcome these deficits by estimating TTC and collision directions from surround-view cameras. As presented in Fig.~\ref{fig:limitations} (b), SCOPE leverages a surround-view camera suite to achieve full 360° environmental perception. This approach immediately overcomes the spatial incompleteness of monocular baselines. Furthermore, we upgrade the prediction target from a scalar TTC value to a vector representation termed \textbf{Oriented TTC (OTTC)}. As shown in Fig.~\ref{fig:limitations} (d), this concept unifies the magnitude of collision urgency (TTC Value) with the relative motion direction (TTC Orientation). By explicitly modeling the direction of collision risk, our method effectively differentiates genuine collision threats from safe bypassing objects, thereby enabling robust and direction-aware collision warning.

To ensure geometric consistency, SCOPE introduces a Range View (RV) representation by projecting multi-view features onto a unified spherical domain. This unification is crucial, as independent predictions on discrete images yield inconsistent results, especially in overlapping regions. Fig.~\ref{fig:range_view} highlights the difference between surround-view representation and RV representation. By establishing a seamless coordinate system, the RV representation bridges the gap between discrete sensor inputs and continuous panoramic perception.

Training such a system requires dense scene flow supervision, yet current public datasets fail to satisfy this requirement. On one hand, the KITTI dataset~\cite{Menze2015CVPR} provides the necessary dense annotations but lacks surround-view coverage. On the other hand, the \textit{nuScenes}~\cite{caesar2020nusc} dataset offers full surround-view images but misses the scene flow data entirely. Although some methods try to generate scene flow labels from nuScenes~\cite{hur2020self, jiang20243dsflabelling}, the results are often too sparse for pixel-level training. To bridge this gap, we reconstruct a densified \textit{nuScenes-SceneFlow} dataset from raw nuScenes point clouds and annotations, generating per-pixel ground truth for OTTC.

% While \textit{nuScenes}~\cite{caesar2020nusc} and Waymo~\cite{sun2020scalability} provide surround-view images, their associated LiDAR point clouds are too sparse to derive dense pixel-level labels. Conversely, the KITTI dataset~\cite{Menze2015CVPR} lacks surround-view coverage. To bridge this gap, we reconstruct a densified \textit{nuScenes-SceneFlow} dataset from raw nuScenes point clouds and annotations, generating per-pixel ground truth for OTTC.

In summary, the main contributions of this paper can be summarized as follows:
\begin{itemize}
  \item [1)] We propose SCOPE, an end-to-end framework for omnidirectional collision prediction. This framework addresses the spatial incompleteness of existing methods, ensuring omnidirectional safety in highly dynamic scenarios.

  \item [2)] We introduce Oriented TTC (OTTC), a novel vector-based risk representation. This formulation resolves the direction ambiguity of traditional TTC, effectively reducing false alarms and accurately ranking collision threats from all directions.

  \item [3)] We design a unified Range View (RV) aggregator. This representation unifies discrete sensor inputs into a continuous domain, ensuring geometric consistency across camera boundaries.

  \item [4)] We construct a densified nuScenes-SceneFlow dataset. This addresses the lack of dense supervision, enabling state-of-the-art (SOTA) performance and zero-shot transfer to a real vehicle platform.
  
\end{itemize}

\section{Related work}
\label{sec:related-work}
In this section, we first survey vision-based collision prediction paradigms, covering BEV object-level pipelines, 3D semantic occupancy, and depth prediction. Then we review TTC estimation methods via optical expansion.

\begin{figure*}[!t]  
    \centering
    \includegraphics[width=1\textwidth]{figs/overview.pdf}
    \caption{Overview of SCOPE. The surround-view images are processed by a backbone to extract multi-scale features. These features are fused across views by the Range View Transformer, transforming the representation from surround view to Range View. The resulting Range View features are then fed into two heads with identical architectures to produce per-pixel OTTC predictions, which is composed by TTC value and TTC orientation.}
    \label{fig:meth}
\end{figure*}

\subsection{Vision-based Collision Prediction Methods}
\label{sec:related-work1}
Existing vision-based collision prediction methods can be broadly grouped into three categories: (i) object detection and tracking–based estimation of time-to-collision, (ii) 3D occupancy–based assessment of free space, and (iii) omnidirectional depth estimation for geometric reconstruction. 

Within the first category, BEV-centric methods are predominant because BEV fuses multi-camera observations in a unified spatial coordinate system with metric-consistent positions and scales, thereby facilitating cross-view and temporal alignment as well as object association. According to the mapping direction from images to BEV, these methods fall into forward projection~\cite{philion2020lss,huang2021bevdet,huang2022bevdet4d,li2023bevdepth,li2023bevstereo,li2024fast} and backward projection~\cite{wang2022detr3d,doll2022spatialdetr,liu2022petr,li2024bevformer}. Forward-projection methods such as LSS~\cite{philion2020lss} model per-ray discrete depth to lift multi-camera features into 3D and then splat to BEV, producing a representation that is robust to depth uncertainty and calibration noise and conducive to cross-camera fusion. Backward projection is represented by BEVFormer~\cite{li2024bevformer}, which employs BEV queries with deformable attention to interact directly with multi-view features and applies temporal attention to historical BEV queries for spatiotemporal alignment. Overall, BEV-based methods effectively unify multi-camera views and enable object-level detection and tracking. However, their reliance on semantic detectors and trackers limits generalization to unseen categories and long-tail objects. 

The second category targets 3D semantic occupancy to infer free space. Most methods predict dense voxel grids of occupied space, enabling collision-risk assessment by measuring spatial proximity~\cite{wei2023surroundocc,yang2024adaptiveocc,li2023fb,huang2023tri,tang2024sparseocc,hou2024fastocc}. For example, SurroundOcc~\cite{wei2023surroundocc} proposes a multi-scale architecture and utilizes multi-scale queries to aggregate 2D features to 3D voxels. Building on SurroundOcc, AdaptiveOcc~\cite{yang2024adaptiveocc} adopts an octree-based representation that uses finer voxels at close range and coarser voxels farther away, substantially reducing compute and memory footprints and enabling in-vehicle deployment. SparseOcc~\cite{tang2024sparseocc} proposes a sparse 3D latent representation that avoids information loss from TPV projections~\cite{huang2023tri}, and alleviates the cubic time-and-memory complexity bottleneck of operating in dense 3D latent spaces. However, unknown or unseen categories are typically handled with a conservative ‘occupied-by-default’ assumption, which often leads to over-braking.

The third category focuses on omnidirectional depth estimation to reconstruct the 3D scene geometry~\cite{wei2023surrounddepth,shi2023ega,kim2022self,schmied2023r3d3,xie2023omnividar,deng2025omnistereo}, thereby supporting collision-risk assessment via geometric analysis. SurroundDepth~\cite{wei2023surrounddepth} employs a cross-view Transformer to interact and fuse features from all cameras, improving cross-view consistency; it also leverages two-view SfM–pretrained sparse pseudo-depth together with extrinsic calibration, allowing the network to learn absolute scale without ground-truth depth. In comparison, EGA-Depth~\cite{shi2023ega} restricts attention to overlapping regions between adjacent views, which reduces computational cost while retaining cross-view coupling. R3D3~\cite{schmied2023r3d3} proposes a multi-camera system for dense 3D reconstruction in dynamic scenes and ego-pose estimation, forming an iterative closed loop between geometric estimation and monocular-depth refinement to obtain consistent, dense reconstructions in complex outdoor environments. In addition, several methods achieve omnidirectional depth with surround fisheye cameras~\cite{xie2023omnividar,deng2025omnistereo}. For example, OmniStereo~\cite{deng2025omnistereo} adopts the Cassini projection to simplify photometric matching under omnidirectional geometry and mitigate distortion, and further designs a lightweight stereo network that attains real-time, full-surround depth while maintaining accuracy.
However, depth encodes only static geometry. In contrast to TTC-based approaches, depth-driven warning methods must additionally estimate velocity or optical flow to yield actionable collision alerts.

\subsection{TTC Estimation Methods}
\label{sec:related-work2}
TTC measures the time to contact along the line of sight. In monocular imagery, it is most reliably inferred from optical expansion, since the rate of image scale growth equals the rate of depth change. Work following this principle broadly falls into two strands: geometry-driven methods and learning-based methods. 

In the first strand, early approaches approximate expansion by bounding-box scale change: a detector provides a temporal sequence of boxes, and changes in box size are mapped to TTC. These approachs inherently limited by category dependence and detectability~\cite{ibrahim2020real,lin2020vision}. Similarly, some methods compute TTC from regional divergence, but at the cost of stronger priors on collision regions and motion patterns~\cite{kilicarslan2018predict}. Other studies derive TTC from optical flow and geometric constraints~\cite{beverley1983texture,camus1995calculating,meyer1994time,swanston1986perceived,schaub2016reactive,yang2020opticalexpansion}. Among them, Schaub et al. provide a concise geometric link between flow and expansion~\cite{schaub2016reactive}, while Yang et al. first regress expansion from flow with a deep network~\cite{yang2020opticalexpansion}. But these methods are sensitive to flow accuracy. To improve robustness, several approaches jointly estimate depth and optical flow~\cite{yin2018geonet,yoon2020novel,luo2019every} or directly infer scene flow~\cite{hur2020self,jiang20243dsflabelling} for TTC.

The second strand employs learning-based scale matching to infer TTC end-to-end, avoiding explicit depth and reducing sensitivity to flow errors. BinaryTTC~\cite{badki2021binaryttc} introduce a scale cost volume, discretizing expansion as scale matching to regress optical expansion; Ling et al. unify expansion and flow within a single model and iteratively refine with a GRU~\cite{ling2022scale, ling2023learning}, achieving higher accuracy at notable computational cost. Advancing this trajectory, FpTTC~\cite{li2024fpttc} treats TTC as a temporal scale matching problem, replacing high-dimensional cost volumes with a multi-level deformable attention scale encoder, and employing a temporal decoder over frames \((t-1,t)\) to predict TTC at time \(t\). This substantially improves both accuracy and speed, making online deployment more practical. However, existing methods are all monocular and forward-view only, and therefore cannot handle lateral or rear collision risks.

To overcome these challenges, we introduce SCOPE, a panoramic framework that augments surround-view perception with orientation-aware TTC estimation. This approach achieves robust, omnidirectional collision warning, freeing the system from the constraints of semantic dependence and limited forward views.

\section{Method}
\label{sec:methd}

This section details our proposed SCOPE method. First, in Section~\ref{sec:prob_form} we introduce the task of panoramic collision prediction and define the concepts of TTC value and TTC orientation, which compose the OTTC. Section~\ref{sec:backbone} presents the backbone of our multi-task prediction network, which enhances features across consecutive frames and models inter-frame correlations. In Section~\ref{sec:rv_agg}, we describe the Range View Transformer that fuses multi-view features via an attention mechanism. Section~\ref{sec:head} details the multi-task prediction head responsible for jointly estimating TTC value and TTC orientation. Finally, Section~\ref{sec:loss} outlines the design of our multi-task loss function.

% We first clarify several related concepts. MiD ($\eta$) refers to the depth ratio between frames, whereas optical expansion, also known as scale, measures the relative change in image size. Mathematically, scale is the reciprocal of MiD ($scale = 1/\eta$). Our model predicts this scale along with orientation, and we subsequently derive TTC from the scale using Eq.~\ref{eq:cam_ttc}.

\subsection{Problem Formulation}
\label{sec:prob_form}

\begin{figure}[!t]  
    \centering
    \includegraphics[width=0.45\textwidth]{figs/optical_expansion_def.pdf}
    \caption{(a) and (b) visualize the relationship between scale change \(s\) and MiD \(\eta\) under scaled orthographic projection. (c) and (d) illustrate the relationship between scale change \(s\) and MiD \(\eta\) when using the spherical projection. According to eq.~\ref{eq:mid_scale} and eq.~\ref{eq:mid_scale_spherical_}, \(s = \theta'/ \theta = l_{sph}'/l_{sph} =  1/\eta\).}
    \label{fig:scale_def}
\end{figure}

We formulate the task of collision prediction as the joint estimation of temporal urgency and spatial motion direction. Specifically, we define the prediction target as an OTTC vector, which consists of a TTC value and a TTC orientation. 

As illustrated in existing TTC prediction methods~\cite{yang2020opticalexpansion, li2024fpttc}, TTC value derived from  Motion-in-Depth (MiD) could be expressed as:

\begin{equation}
\label{eq:cam_ttc}
    \mathrm{TTC}
    = Z/{\frac{dZ}{dt}}
    = Z/{\frac{Z - Z'}{\Delta t}}
    = \frac{\Delta t}{1 - \eta},
    \text{where }\eta \coloneqq \frac{Z'}{Z}.
\end{equation}
where \(Z\) and \(Z'\) denote the pixel-wise depths at adjacent times, respectively. The interval between frames is \(\Delta t\), and the MiD \(\eta\) represents the per-pixel depth ratio between consecutive frames. Fig.~\ref{fig:scale_def} (a) and (b) show the relationship between MiD and the scale change of a planar object in the image plane. If we assume that the moving object is rigid and the camera model is scaled orthographically, the MiD \(\eta\) is inversely proportional to the scale change \(s\) in the image plane:
\begin{equation}
\label{eq:mid_scale}
    \frac{l}{L}=\frac{f}{Z}, \frac{l'}{L}=\frac{f}{Z'} \quad \Rightarrow \quad \eta = \frac{Z'}{Z} = \frac{l}{l'} = \frac{1}{s}
\end{equation}
where \(l\) and \(l'\) denote the lengths of the object in the image plane at adjacent timestamps, \(L\) is the actual length and \(f\) is the focal length of the camera. Therefore, TTC can be directly derived from Eq.~\ref{eq:cam_ttc} by simply estimating the inter-frame scale change \(s\).

In this paper, we expand the existing forward view TTC prediction to panoramic TTC prediction by adopting a spherical projection. Rather than projecting the scene geometry onto a planar image, we project it onto the surface of a sphere centered at the surround-view camera system, establishing a unified RV representation. An illustration of the spherical projection geometry and the resulting TTC–MiD relationship is shown in Fig.~\ref{fig:scale_def} (c) and (d). The scale change of the object in the spherical projection surface can be calculated as follows:
\begin{equation}
\label{eq:mid_scale_spherical}
    \eta = \frac{l_{sph}}{l_{sph}'} = \frac{\theta}{\theta'} = \frac{\arctan(\frac{L}{2Z})}{\arctan(\frac{L}{2Z'})}
\end{equation}
where \(l_{sph}\) and \(l_{sph}'\) denote the lengths of the object in the spherical projection surface at adjacent times, respectively. \(\theta\) and \(\theta'\) denote the corresponding angles. Since our method operates at a pixel level, \(L\) represents the physical spatial extent covered by a single pixel's field of view. Given that this length is negligible compared to the depth (\(L \ll Z\)) in typical driving scenarios, we can apply the small-angle approximation:
\begin{equation}
\label{eq:mid_scale_spherical_}
    \arctan(\frac{L}{2Z}) \approx \frac{L}{2Z}\quad \Rightarrow \quad \frac{l_{sph}}{l_{sph}'} \approx \frac{Z'}{Z} = \eta
\end{equation}
which means that the MiD can still be approximated by the scale change of the object in the spherical projection surface. Therefore, we can still calculate TTC in the spherical projection surface by estimating scale changes under small-angle approximation and typical driving distances.

Complementing the TTC value, we define TTC orientation to characterize the spatial direction of collision risk. Specifically, it is computed as the horizontal angle between the relative motion vector and the inward radial direction.Concretely, given two consecutive sweeps at times \(t\) and \(t{+}\Delta t\), let \(\mathbf{p}_i=(x_i,y_i,z_i)\) and \(\mathbf{p}'_i=(x'_i,y'_i,z'_i)\) denote the 3D coordinates of the \(i\)-th point correspondence, and let \(\mathbf{p}^{xy}_i=[x_i\;y_i]^\top\) and \(\mathbf{p}'^{\,xy}_i=[x'_i\;y'_i]^\top\) be their horizontal components, as vertical motion is ignored by design. We form the horizontal motion vector \(\mathbf{m}_i=\mathbf{p}'^{\,xy}_i-\mathbf{p}^{xy}_i\) and the inward radial vector \(\mathbf{u}_i=-\,\mathbf{p}^{xy}_i\). The orientation \(\theta_i\in[0,\pi]\) is then computed as
\begin{equation}
\label{eq:orientation}
    \theta_i \;=\; \arccos(\frac{\mathbf{m}_i^\top \mathbf{u}_i}{\lVert\mathbf{m}_i\rVert_2 \lVert\mathbf{u}_i\rVert_2})
\end{equation}
Intuitively, \(\theta_i\approx 0\) indicates the object is approaching, \(\theta_i\approx \pi\) means it is moving away, and \(\theta_i=\tfrac{\pi}{2}\) represents a passing-by motion. We filtered out invalid pixels as well as points with negligible horizontal motion (i.e., \(\lVert\mathbf{m}_i\rVert_2 < \epsilon\)) to avoid undefined orientations.

Finally, we integrate the temporal urgency and spatial direction into the unified OTTC vector. For each pixel \(i\), the OTTC vector \(\mathbf{v}_{i}\) is defined by combining the TTC value derived from Eq.~\ref{eq:cam_ttc} and the TTC orientation \(\theta_i\) from Eq.~\ref{eq:orientation}:
\begin{equation}
\label{eq:ottc_vector}
\mathrm{OTTC}_i = [\mathrm{TTC}_i, \theta_i]^\top
\end{equation}
This vector representation serves as the comprehensive prediction target of our framework, enabling the distinction between genuine collision threats and safe dynamic objects. 

Detailed information on constructing the pixel-wise TTC value and TTC orientation maps from LiDAR point clouds, including the spherical projection procedure, is available in Section \ref{sec:rv_gt}.

\subsection{Network Backbone}
\label{sec:backbone}
Guided by the problem formulation, our network is designed to jointly regress pixel-wise TTC value and TTC orientation. Since both scale change and orientation arise from frame‐to‐frame differences, they can be recovered via feature matching like optical flow~\cite{xu2022gmflow}. 
To realize this, we propose a three‐stage architecture comprising attention enhancement, correlation inference, and feature guided propagation, thereby enabling accurate omnidirectional collision prediction. Given multi-level features \(F_{t-1}^{(l)}, F_{t}^{(l)} \in \mathbb{R}^{C\times H_l\times W_l}\), where \(l\in \left\{ 0,1,\dots,L-1 \right\} \) is the pyramid feature level, \(H_l\) and \(W_l\) are the height and width, respectively, and \(C\) is the number of channels. We first reorganize and align the two frames by self-attention and cross-attention modules based on swin-transformer~\cite{liu2021swin}, propagating intra-frame structure and inter-frame alignment, yielding enhanced features \(\hat{F}_{t-1}^{(l)}\) and \(\hat{F}_{t}^{(l)}\). On top of the enhanced features, we infer a geometric prior via soft matching, which is implemented by a correlation layer~\cite{xu2022gmflow}. Let \(\Omega_l = \left\{1,\dots,H_l \right\} \times \left\{1,\dots,W_l \right\}\) denote the pixel grid at level \(l\), and \(\mathbf{p}_i^{(l)}\in \mathbb{R}^2\) be the 2D coordinate of the \(i\)-th pixel in \(\Omega_l\), \(\hat f_{t-1,i}^{(l)},\hat f_{t,i}^{(l)} \in \mathbb{R}^C\) be the corresponding features in \(\hat{F}_{t-1}^{(l)}\) and \(\hat{F}_{t}^{(l)}\), respectively. For a source position \(\mathbf{p}_i^{(l)}\), calculate similarity and matching distribution as:

\begin{equation}
\label{eq:sim_and_match}
    S_{ij}^{(l)}=\dfrac{\langle \hat f_{t-1,i}^{(l)},\,\hat f_{t,j}^{(l)}\rangle}{\sqrt{C}},\quad \pi_{ij}^{(l)}=\dfrac{e^{S_{ij}^{(l)}}}{\sum_{k} e^{S_{ik}^{(l)}}}
\end{equation}
The destination position \(\mathbf{c}_i^{(l)}=\sum_j\pi_{ij}^{l}\mathbf{p}_j^{(l)}\), and the initial optical flow feature is: 
\begin{equation}
\label{eq:flow_feat}
    \mathbf{u}_i^{(l)} = \mathbf{c}_i^{(l)} - \mathbf{p}_i^{(l)} \in \mathbb{R}^2
\end{equation}
We again employ an attention mechanism to update the optical flow feature on same feature level. Specifically, we compute similarities using the feature map \(\hat{F}_{t-1}^l\) and \(\hat{F}_{t}^l\) as both query and key, and then apply the resulting attention weights to refine the optical flow feature, thereby achieving feature-guided propagation, yielding refined optical flow feature \(\mathbf{\tilde{u}}_i^{(l)}\). If \(l<L-1\), we upsample the optical flow feature \(\mathbf{\tilde{u}}^{(l)}\) to the next finer level \(l+1\) using bilinear interpolation:
\begin{equation}
\label{eq:upsample}
    \mathbf{\tilde{u}}_{\uparrow}^{(l \rightarrow l+1)} = \operatorname{BilinearUp}(\mathbf{\tilde{u}}^{(l)},\, \text{scale}=2)\times 2
\end{equation}

Following these three stages, we obtain multi-level enhanced features \(\hat{F}_{t-1}^{(l)}\) and \(\hat{F}_{t}^{(l)}\), as well as the initial optical flow feature \(U\) from the coarsest level. We also use the optical flow feature \(U\) to initialize the scale and orientation prediction head as FP-TTC~\cite{li2024fpttc}. 

\begin{figure*}[!t]  
    \centering
    \includegraphics[width=0.98\textwidth]{figs/Dataset_gen_pipeline.pdf}
    \caption{The pipeline for generating range view ground truth from \textit{nuScenes} dataset. We first build a dense static map and per-instance dynamic templates by aggregating multi-frame point clouds. For each pair of consecutive frames, we then obtain matched point sets by intersecting their point clouds. Finally, we project the point correspondences onto the range view grid to generate per-pixel TTC Value and TTC Orientation maps.}
    \label{fig:dataset_gen_pipeline}
\end{figure*}

\subsection{Range View Transformer}
\label{sec:rv_agg}
The Range View Transformer (RVT) performs geometrically consistant, learnable fusion of multi-view features on a range view grid. We use this module to aggregate the multi-view features \(F_{t-1}\), \(F_{t}\) and the initial optical flow feature \(U\) into range view. Multi-scale deformable attention (MSDA)~\cite{zhu2020deformable} is applied to implement RVT. 

Let the target range view grid be \((H_r, W_r)\), and each cell \((i,j)\) defines a ray direction \(\mathbf{r}_{ij}\). Sampling \(K\) depths \(\{d_k\}_{k=1}^K\) along this ray yields 3D points \(\mathbf{X}_{ij}^{(k)} = d_k\,\mathbf{r}_{ij}\). For each camera \(\ell \in \{1,\dots,L\}\), 3D points could be projected to image space with calibrated intrinsic and extrinsic: 
\begin{equation}
\label{eq:ref_points}
\mathbf{u}_{ij\ell}^{(k)} = \pi(\mathbf{K}_\ell [\,\mathbf{R}_{\ell}\,|\,\mathbf{t}_{\ell}\,]\, \mathbf{X}_{ij}^{(k)})
\end{equation}
where \(\pi\) denotes perspective division, \(K_\ell\) and \([\mathbf{R}_{\ell}\,|\,\mathbf{t}_{\ell}]\) are the intrinsic and extrinsic of camera \(\ell\), respectively. The resulting \(\mathbf{u}_{ij\ell}^{(k)}\) are normalized to \([0,1]^2\) and accompanied by a visibility mask \(m_{ij\ell}^{(k)}\) encoding occlusion and field-of-view validity. Based on this geometry, RVT supplies explicit geometry to deformable attention by forming sampling locations \(\mathbf{L}_{ij} = \{\mathbf{u}_{ij\ell}^{(k)}\} \in [0,1]^2\) and a mask \(m_{ij\ell}^{(k)}\). 

Let the feature from the camera \(\ell\) be \(F_{\ell}\in \mathbb{R}^{C\times H_{\ell} \times W_{\ell}}\) . We flatten along the spatial axes and concatenate across cameras along the sequence dimension, yielding the value:
\begin{equation}
\label{eq:rva_value}
\mathbf{V} = \left [X_{1}, X_{2}, \dots, X_{L} \right]\in \mathbb{R}^{S\times C},\quad  S=\sum_{\ell=1}^{L}H_{\ell}W_{\ell}
\end{equation}
while the query feature could be derived as follows:
\begin{equation}
\label{eq:rva_query}
\mathbf{Q} = \operatorname{Flatten}\bigl(\phi{(\tilde{Q})} + \operatorname{PE}(\phi(\tilde{Q}))\bigr) \in \mathbb{R}^{(H_rW_r)\times C}
\end{equation}
where \(\phi\) is a light-weight convolution block, and \(\tilde {Q} \in \mathbb{R}^{C\times H_r\times W_r}\) is the initial query feature. The query-derived weights \(\tilde{\alpha}_{ij\ell}^{(k)} = \mathrm{softmax}(\phi(\mathbf{Q}_{ij}))\) are combined with visible mask and renormalized, thereby we restrict attention to visible locations. 
% as follows:
% \begin{equation}
% \label{eq:rva_renorm}
% \alpha = \frac{\tilde{\alpha}_{ij\ell}^{(k)}m_{ij\ell}^{(k)}}{\sum_{\ell',k'}\tilde{\alpha}_{ij\ell'}^{(k')}m_{ij\ell'}^{(k')}}
% \end{equation}
% thereby we restrict attention to visible locations. 

Applying MSDA separately at both times steps \(t-1\) and \(t\) with inputs  \((\mathbf{V},\mathbf{L},\alpha)\) performs cross-camera, cross-depth soft aggregation and yields fused range view features \(F_{t-1, r}, F_{t, r}  \in \mathbb{R}^{H_r \times W_r \times C}\).  If we use optical flow feature \(U\) to construct the value feature, we will get the range view features \(U_r\). Compared to hard projection, RVT avoids single-pixel indexing and single-depth assumptions that cause seams and occlusion artifacts; compared to purely learned alignment, it reduces the search space and improves stability and interpretability. The fused representation \(F_r, U_r\) serves as a cross-view consistent, panoramic feature for downstream tasks such as scale estimation and orientation prediction.

To stabilize early training of the range view module, we optionally employ a monocular depth network, DepthAnything (DA)~\cite{depth_anything_v2}, as a teacher. DA produces per-camera depth maps that are projected into the spherical range view grid and used as a coarse geometric prior. During training, this prior serves two roles: (i) it initializes the query features for our range view aggregation, and (ii) it provides an auxiliary distillation signal on the aggregated range view features. The weight of the distillation loss is gradually annealed to zero so that, at convergence, the model no longer depends on DA predictions and performs inference purely from surround-view images. All quantitative results reported in Sec.~\ref{sec:exp} are obtained with the final teacher-free model, while Sec.~\ref{sec:ablation} further studies alternative query initialization strategies with and without teacher guidance.


\subsection{TTC Value and TTC Orientation Head}
\label{sec:head}
We use two prediction heads to jointly estimate the TTC value and TTC orientation on the range view grid. Both the TTC value head and the TTC orientation head share a unified pipeline, which consists of an MSDA-based encoder and a RAFT-based recurrent refinement decoder.

The encoder has similar architecture as the RVT module, and the range view optical flow feature \(U_r\) is applied to get the initialized query feature \(\tilde{Q}\) for the TTC value and orientation heads. The query features and value features are constructed following Eq.~\ref{eq:rva_query} and Eq.~\ref{eq:rva_value}. Unlike RVT, this encoder does not rely on additional geometric information to locate reference points, because the queries and values have already been aligned within the RVT module. Consequently, it only uses a lightweight convolutional block to predict the sampling offsets \(\operatorname{Bias}=\phi(\mathbf{Q}) \in \mathbb{R}^{H_r\times W_r \times 2}\). 

After the MSDA encoder consumes \((\mathbf{Q},\mathbf{V},\mathbf{L},\alpha)\) and yields fused range view features \(F_{t-1,r},\,F_{t,r}\), and the range view optical flow \(U_r\), the decoder refines a coarse initialization \(s^{(0)}\) on the \((H_r,W_r)\) grid into a high-resolution map. The coarse seed can be regressed from \([\,U_r,\,F_{t,r}\,]\). The decoder input feature is
\begin{equation}
\label{eq:decoder_input}
E_r = \phi\ \big([\;F_{t-1,r},\,F_{t,r}\;]\big)
\end{equation}

where \(\phi(\cdot)\) denotes lightweight convolutional fusion. A correlation-derived guidance cue is denoted by \(A_r\), providing dynamic evidence during refinement. The decoder first merges \([\,E_r,\,F_{t,r}\,]\) into a static context \(z\)  and initializes the hidden state \(n^{(0)}\). At refinement step \(\tau\) (\(\tau=1,\dots,T\)), a motion encoder produces
\begin{equation}
\label{eq:motion_encoder}
m^{(\tau)} \;=\; \mathcal{E} \big(s^{(\tau-1)},\,A_r\big)
\end{equation}
after which a recurrent unit updates the state \(n^{(\tau)}\) from \([\,z,\,m^{(\tau)}\,]\). A shallow prediction head yields the increment \(\Delta s^{(\tau)}\), and the estimate is accumulated residually:
\begin{equation}
\label{eq:scale_accum}
s^{(\tau)} \;=\; s^{(\tau-1)} \;+\; \Delta s^{(\tau)}
\end{equation}
After \(T\) iterations, the low-resolution output \(s^{(T)}\) is upsampled via RAFT-style convex combination: weights predicted from \([\,s^{(T)},\,F_{t,r}\,]\) blend the \(3\times3\) neighborhood \(\mathcal{N}(\mathbf{p})\) to reconstruct a high-resolution location \(\mathbf{p}+\boldsymbol{\delta}\) :
\begin{equation}
\hat{s}(\mathbf{p}+\boldsymbol{\delta}) \;=\;
\sum_{\mathbf{q}\in\mathcal{N}(\mathbf{p})}
w_{\mathbf{q}}(\mathbf{p},\boldsymbol{\delta})\,s^{(T)}(\mathbf{q}),
\quad
\sum_{\mathbf{q}} w_{\mathbf{q}}(\mathbf{p},\boldsymbol{\delta}) \;=\; 1
\end{equation}
In this way, the decoder couples coarse-to-fine recurrent refinement with learnable convex upsampling that avoids seams and artifacts of hard indexing. 

\begin{figure*}[!t]  
    \centering
    \includegraphics[width=0.98\textwidth]{figs/pc_overlay.pdf}
    \caption{Comparison of single-frame and multi-frame ground-truth generation. By aggregating multiple frames, the point cloud becomes denser and more complete, capturing richer texture and edge information, which facilitates more accurate per-point annotations of TTC value and orientation.}
    \label{fig:pc_overlay}
\end{figure*}

\subsection{Training Loss}
\label{sec:loss}
We evaluate errors only on supervised pixels indicated by a binary mask $m_{ij}\in\{0,1\}$ by defining the index set \(\Omega=\{(i,j)\mid m_{ij}=1\}\).
For TTC value branch, we parameterize the prediction using a positive scale change $s_{ij}>0$ with positive ground truth $\hat s_{ij}>0$. We then utilize a log-domain discrepancy that emphasizes relative error:
\begin{equation}
\label{eq:scale_loss}
\mathcal{L}_{\text{value}}
=\frac{1}{|\Omega|}\sum_{(i,j)\in\Omega}
\bigl|\log(s_{ij})-\log(\hat s_{ij})\bigr|
\end{equation}
For orientation prediction $\theta_{ij}\in[0,\pi]$ with ground truth
$\hat\theta_{ij}\in[0,\pi]$, we adopt a mean-squared error to provide a smooth,
convex objective over the bounded interval:
\begin{equation}
\label{eq:orientation_loss}
\mathcal{L}_{\text{ori}}
=\frac{1}{|\Omega|}\sum_{(i,j)\in\Omega}\bigl(\theta_{ij}-\hat\theta_{ij}\bigr)^2
\end{equation}

If $|\Omega|=0$, both losses are defined as zero to avoid penalizing unsupervised regions. Since both branches are trained jointly, the overall loss is a weighted sum:
\begin{equation}
\label{eq:total_loss}
\mathcal{L}_{\text{total}}
=\lambda_{\text{value}}\mathcal{L}_{\text{value}} + \lambda_{\text{ori}}\mathcal{L}_{\text{ori}}
\end{equation}
where we set \(\lambda_{\text{value}}=1\) and \(\lambda_{\text{ori}}=1\) in all experiments.

\section{Range View Ground Truth Generation}
\label{sec:rv_gt}
To generate ground-truth TTC value and orientation, we leverage scene flow, which provides per-point 3D motion. Using the established kinematic equations Eq.~\ref{eq:mid_scale_spherical_} and Eq.~\ref{eq:orientation}, these motions could be converted into per-point TTC value and orientation labels. Given that our method performs vision-based panoramic perception, we adopt the \textit{nuScenes}~\cite{caesar2020nusc} dataset, which offers 360° surround view images together with synchronized LiDAR point clouds, making it therefore well-suited for deriving and validating the required annotations. The range view ground truth is constructed based on these two steps: (i) multi-frame point cloud accumulation to densify the sampling of the scene. Fig.~\ref{fig:pc_overlay} presents a comparison between sparse and dense point clouds, where detailed textures can be clearly observed in the densified results. And (ii) projection of the accumulated point cloud onto the range view grid to generate per-pixel TTC value and TTC orientation maps.

\subsection{Multi-Frame Point Cloud Overlay}
\label{sec:multi-frame}

We construct a temporally aggregated static map together with per-instance dynamic templates~\cite{wei2023surroundocc} for each \textit{nuScenes} scene. For each pair of consecutive frames, we then obtain matched point sets by intersecting their point clouds. The resulting point-wise correspondences form a dense scene flow field, which provides the basis for generating TTC value and TTC orientation ground truth. The overall pipeline is illustrated in Fig.~\ref{fig:dataset_gen_pipeline}. 

We process each frame \(t\) in a scene. By using the annotated 3D boxes, we remove both the in-box points and the ego-near points to obtain a static point cloud subset \(\mathcal{S}_t\). With the provided ego-poses and extrinsic calibration, we compute a rigid transform \(\mathbf{H}_{t\rightarrow 0}\) that maps each LiDAR frame \(L_t\) to a common reference frame \(L_0\). All static subsets are then warped into \(L_0\) and united over time, yielding a temporally aggregated dense static map:
\begin{equation}
\label{eq:static_map_points}
\mathcal{S}^{(\mathrm{ref})} = \bigcup_{t} \mathbf{H}_{t\rightarrow 0}\!\bigl(\mathcal{S}_t\bigr)
\end{equation}
For dynamic objects, we traverse all frames in the scene and group points by instance token \(q\). In each frame where instance \(q\) appears, we extract the in-box points and normalize them to an instance-centric local frame (translation to the box center and rotation to cancel the box heading), yielding a per-frame object segment \(\mathcal{O}_{q,t}\). Aggregating these object segments over time produces a dense, multi-view object segment for instance \(q\): 
\begin{equation}
\label{eq:object_segment_points}
\mathcal{O}_{q}=\bigcup_{t}\mathcal{O}_{q,t}
\end{equation}

To synthesize one-to-one corresponding point sets for consecutive frames \((t{-}1,t)\), we first project the aggregated static map from the reference frame back into each LiDAR frame via \(\mathbf{H}_{0\rightarrow t-1}\) and \(\mathbf{H}_{0\rightarrow t}\), obtaining \(\mathcal{S}^{L_{t-1}}\) and \(\mathcal{S}^{L_{t}}\). For dynamic objects, we identify instance tokens that appear in both frames,
\(\mathcal{Q}_{t-1,t}\). For each \(q \in \mathcal{Q}_{t-1,t}\), we retrieve its
object segment \(\mathcal{O}_{q}\) and place it into each frame according to the
annotated object poses and ego-poses, yielding dynamic point sets
\begin{equation}
\label{eq:dynamic_point_sets}
\mathcal{D}^{L_{t-1}} = \bigcup_{q \in \mathcal{Q}_{t-1,t}} \mathcal{O}_{q}^{L_{t-1}}, 
\quad
\mathcal{D}^{L_{t}} = \bigcup_{q \in \mathcal{Q}_{t-1,t}} \mathcal{O}_{q}^{L_{t}}
\end{equation}
The final paired point sets
\(\mathcal{P}^{t-1}=\mathcal{S}^{L_{t-1}}\cup\mathcal{D}^{L_{t-1}}\) and
\(\mathcal{P}^{t}=\mathcal{S}^{L_{t}}\cup\mathcal{D}^{L_{t}}\) are then cropped to a shared axis-aligned region.
Grounded in the rigid mappings \(\mathbf{H}_{t\rightarrow 0}\) and \(\mathbf{H}_{0\rightarrow t}\), this pipeline enforces cross-frame geometric consistency, augments each frame with a dense static context and stable per-instance geometry, and produces clean, temporally aligned point correspondences that underpin scene flow estimation and subsequent scale/orientation supervision.



\subsection{Project Accumulated Point Cloud to Range View}
\label{sec:rv_project}
Given the accumulated and consecutive point sets \(\mathcal{P}^{t-1}\) and \(\mathcal{P}^{t}\), we construct TTC value and TTC orientation ground truth map as follows. For each point \(i\), we first obtain the planar depths \(d^{t-1}_i\) and \(d^{t}_i\) from their \(xy\)-norms at times \(t-1\) and \(t\), respectively. The scale is then computed as the ratio \(s_i = d^{t}_i / d^{t-1}_i\) based on Eq.~\ref{eq:mid_scale_spherical_}, and then TTC value is computed via Eq.~\ref{eq:cam_ttc}. 
We next form the \(xy\)-plane motion vector \(\mathbf{m}_i\) between \({\mathbf{p}}^{t-1}_i\) and \(\mathbf{p}^{t}_i\), and the inward radial unit vector \(\hat{\mathbf{u}}_i\) pointing from \({\mathbf{p}}^{t-1}_i\) toward the LiDAR origin. The orientation \(\theta_i \in [0,\pi]\) is defined as the angle between \(\mathbf{m}_i\) and \(\hat{\mathbf{u}}_i\) according to Eq.~\ref{eq:orientation}.
% Because camera captures in \textit{nuScenes} are triggered by the LiDAR, each image frame is strictly synchronized with a corresponding LiDAR sweep and thus has a one-to-one associated point set \(\mathbf{P}_t\). Consequently, when acquiring an image group \(\mathcal{I}_t\), we simultaneously obtain the matched LiDAR point set \(\mathbf{P}_t\). 

% For each point we compute planar depths \(d^{t-1}_i\) and \(d^{t}_i\) from their \(xy\)-norms; the scale is the ratio \(s_i = d^{t}_i / d^{t-1}_i\) based on Eq.~\ref{eq:mid_scale_spherical_}. We form the \(xy\)-plane motion vector \(\mathbf{m}_i\) between \({\mathbf{p}}^{t-1}_i\) and \(\mathbf{p}^{t}_i\), and the inward radial unit direction \(\hat{\mathbf{u}}_i\) pointing from \({\mathbf{p}}^{t-1}_i\) toward the LiDAR origin; the orientation \(\theta_i \in [0,\pi]\) is defined as the angle between \(\mathbf{m}_i\) and \(\hat{\mathbf{u}}_i\). 

We map 3D points to a range view via a spherical projection \(\Pi:\ \mathbb{R}^3 \longmapsto \mathbb{R}^2\) that converts each point \(\mathbf{p}_i=(x_i,y_i,z_i)\) into pixel coordinates \((u_i,v_i)\) on an image of size \(H\times W\)~\cite{milioto2019rangenet++}. The mapping uses the yaw \(\psi_i\) and pitch \(\phi_i\) derived from the point’s direction and then normalizes them to the discrete pixel domain, as written below:

\begin{equation}
    \begin{aligned}
    u_i &= \tfrac{1}{2}\!\left(\psi_i\,\pi^{-1} + 1\right)\,W\\
    v_i &= \left(1 - (\phi_i + f_{\text{down}})\,f^{-1}\right)\,H
    \end{aligned}
    \label{equ:proj}
\end{equation}
    
where \(\psi_i = -\arctan(y_i, x_i)\), \(\phi_i = \arcsin(z_i / \lVert \mathbf{p}_i\rVert_2)\), \(f = f_{up} + f_{down}\) is the vertical field of view (FOV). 
During projection, points are depth-sorted and rasterized so that nearer points overwrite farther ones. 
This back-to-front procedure generates three per-pixel maps for both frames \(t-1\) and \(t\): scale \(\mathbf{S}\in\mathbb{R}^{H\times W}\) (used to derive \textbf{TTC value}) with entries \(S_{v_i,u_i}=s_i\), orientation \(\mathbf{O}\in\mathbb{R}^{H\times W}\) with \(O_{v_i,u_i}=\theta_i\), and a validity mask \(\mathbf{M}\in\{0,1\}^{H\times W}\) indicating whether a pixel has a corresponding point.

\section{Experiment}
\label{sec:exp}

\begin{figure*}[t]
    \centering
    \begin{minipage}{0.98\textwidth}
        \centering
        \includegraphics[width=\linewidth]{figs/nusc_result0.pdf}
    \end{minipage}

    \begin{minipage}{\textwidth}
        \centering
        \includegraphics[width=\linewidth]{figs/colorbar.pdf}
    \end{minipage}
    \caption{\textbf{Scale} and \textbf{orientation} prediction result in the range view domain on the \textit{nuScenes-SceneFlow} test split. It shows the predicted scale and orientation maps on the spherical range view, where approaching motion along the viewing rays is clearly highlighted. \protect \legendbox{darkred}: In this scenario, a fast-approaching vehicle from the rear direction is accurately emphasized as high-risk regions in the predicted fields. The left color bar corresponds to the \textbf{TTC Value}, where \(\text{TTC Value} < 0\) indicates objects moving away from the ego vehicle with no collision risk. The right color bar corresponds to the \textbf{TTC Orientation}.}
    \label{fig:nusc_result}
\end{figure*}

% \begin{table}[t]
% \centering
% \caption{Evaluation of MiD and Percentage Errors on nuScenes-SceneFlow}
% \label{table:mid_err}
% \setlength{\tabcolsep}{4pt} 
% \footnotesize               
% \begin{adjustbox}{width=\columnwidth,center}
% \begin{tabular}{@{}l l c c c c c@{}}
% \toprule
% \textbf{Input} & \textbf{Method} & \textbf{MiD Err} (\(\downarrow\)) & \textbf{Err-1} (\(\downarrow\)) & \textbf{Err-2} (\(\downarrow\)) & \textbf{Err-5} (\(\downarrow\)) & \textbf{latency(ms)} \\
% \midrule
% \multirow{4}{*}{Mono}
%     & Expansion~\cite{yang2020opticalexpansion} & 79.26           & 0.08          & 15.22         & 18.95         & 67.53 \\
%     & BiTTC~\cite{badki2021binaryttc}           & 75.37           & 0.03          & 10.16         & 16.45         & 76.34 \\
%     & FP-TTC~\cite{li2024fpttc}                  & 58.68           & 0.04          & 7.84          & 12.69        & \textbf{35.40} \\
%     & \textbf{Ours}                             & \textbf{43.23}  & \textbf{0.01} & \textbf{0.47} & \textbf{0.77} & -- \\
% \midrule
% \multirow{4}{*}{Pano}
%     & Expansion~\cite{yang2020opticalexpansion} & 98.04          & 0.33          & 7.02          & 10.19         & 392.40 \\
%     & BiTTC~\cite{badki2021binaryttc}           & 86.58          & \textbf{0.10} & 5.43          & 8.82          & 461.60 \\
%     & FP-TTC~\cite{li2024fpttc}                  & 70.03          & 0.11          & 5.07          & 8.47          & 255.66 \\
%     & \textbf{Ours}                             & \textbf{57.42} & 0.25          & \textbf{2.04} & \textbf{3.33} & \textbf{102.32} \\
% \bottomrule
% \end{tabular}
% \end{adjustbox}
% \end{table}

% \begin{table}[t]
% \centering
% \caption{Frame-level Recall and False Alarm Rate (FAR) on Local Platform}
% \label{table:local_recall_far}
% \setlength{\tabcolsep}{3pt} 
% \small 
% \begin{adjustbox}{max width=\columnwidth,center}
% \begin{tabular}{@{}l l c c@{}}
% \toprule
% \textbf{Input} & \textbf{Method} & \textbf{Recall} (\(\uparrow\)) & \textbf{FAR} (\(\downarrow\)) \\
% \midrule
% \multirow{4}{*}{Pano}
% &Expansion~\cite{yang2020opticalexpansion} & 0.863          & 0.063   \\
% &BiTTC~\cite{badki2021binaryttc}           & 0.917          & 0.031   \\
% &FP-TTC~\cite{li2024fpttc}                 & 0.955          & 0.020   \\
% &\textbf{Ours}                             & \textbf{0.976} & \textbf{0.011} \\
% \bottomrule
% \end{tabular}
% \end{adjustbox}

% \vspace{2pt}
% {\footnotesize
% \textit{Note:} Recall is computed over frames with ground-truth collision risk,
% while FAR is computed over frames without ground-truth collision risk.
% }
% \end{table}
    
    


\subsection{Datasets and Real Vehicle Platform Setup}

Using our proprietary \textit{nuScenes-SceneFlow} dataset and the procedure in Sec.~\ref{sec:rv_gt}, we derive range view, per-pixel TTC value and orientation ground truths for 20 \textit{nuScenes} scenes, yielding 2{,}155 pairs of surround-view images with corresponding annotations. The temporal interval between the two images in each pair is 0.1s. We adopt a scene-level split of 70\%/10\%/20\% for training, validation, and testing, ensuring that samples from the same scene do not appear across different splits. 

Qualitative results on the test split are shown in Fig.~\ref{fig:nusc_result}. 
From top to bottom, each example displays the RGB image, the predicted and ground-truth TTC value maps, and the predicted and ground-truth TTC orientation maps in the unified range view domain. 
The predicted TTC value map closely follows the ground truth. On the TTC value map, warm colors concentrating on objects that rapidly approach the ego vehicle along the viewing rays, and cooler colors on static background regions, indicating large or negative TTC with no collision risk. 
The rear approaching vehicle highlighted in Fig.~\ref{fig:nusc_result} is accurately emphasized as a high-risk region in the predicted scale map, while surrounding area is correctly suppressed as low-risk. 
Similarly, the predicted TTC orientation map aligns well with the ground truth, capturing the direction of relative motion for moving vehicles and remaining stable on static structures. 
These observations confirm that SCOPE learns a geometrically consistent OTTC field, which jointly encodes collision urgency and motion direction across all directions.

The hardware configuration of the real vehicle platform is shown in Fig.~\ref{fig:hareware}. It uses six surround-view cameras of type SG2-IMX390C-5200-GMSL-H100F1, each with a horizontal field of view of 100° and a vertical field of view of 53°.
Additionally, our end-to-end PyTorch implementation achieves a processing speed of approximately \(10\,\mathrm{Hz}\) on a single NVIDIA RTX 4090D, satisfying typical real-time latency requirements for onboard collision warning, motion planning, and control in autonomous driving systems. 

% In addition, we collected surround-view data in campus environments to evaluate cross-domain robustness and generalization. The campus set comprises two collision-free free-driving sequences (8{,}730 surround-view image pairs) and eight real-collision sequences (660 image pairs). The temporal interval between the images in each pair is 0.1,s. The prediction results on the real-collision scenarios are illustrated in Fig.~\ref{fig:campus_zeroshot}.

\subsection{Training Setup}

We jointly train the TTC value and orientation heads using the range view ground truths derived from the \textit{nuScenes-SceneFlow} data. Training is performed on \(4\times\) NVIDIA A800 GPUs for a total of 94{,}000 iterations. We use AdamW as the optimizer with an initial learning rate of \(4\times 10^{-5}\), and adopt OneCycleLR as the learning-rate scheduling policy. Unless otherwise specified, all testing is carried out on a single NVIDIA GeForce RTX 4090D.

\subsection{Evaluating Metrics}

\subsubsection{MiD Error}
We use the metric MiD error to assess the accuracy of the TTC value prediction results~\cite{li2024fpttc}. MiD reflects the scale change in the depth direction during the motion, which can be computed from the depth ratio as in Eq.~\ref{eq:mid_scale_spherical_}. The MiD error is defined as
\begin{equation}
\label{eq:mid_err}
\mathrm{MiD}_{\mathrm{err}}=\bigl|\log(\eta_{\mathrm{gt}})-\log(\eta)\bigr|\times 10^{4}
\end{equation}

Beyond MiD, we further evaluate \(\mathrm{TTC}\) as a binary decision under safety-relevant time thresholds. Over the valid pixel set \(\Omega\), for a threshold \(T\in\{1,2,5\}\) seconds we define the ground-truth and predicted binary labels
\begin{equation}
\label{eq:ttc_bin}
y^{\mathrm{gt}}_{T}(p)=\mathbbm{1}\!\big(\tau_{\mathrm{gt}}(p)<T\big),\qquad
\hat{y}_{T}(p)=\mathbbm{1}\!\big(\tau(p)<T\big)
\end{equation}
where \(\tau_{\mathrm{gt}}(p)\) and \(\tau(p)\) denote the ground-truth and predicted TTC, respectively. The error rate at threshold \(T\) is
\begin{equation}
\label{eq:err_T}
\mathrm{Err}\text{-}T=\frac{1}{|\Omega|}\sum_{p\in\Omega}\mathbbm{1}\!\big(\hat{y}_{T}(p)\neq y^{\mathrm{gt}}_{T}(p)\big)
\end{equation}
We report \(\mathrm{Err}\text{-}1\), \(\mathrm{Err}\text{-}2\), and \(\mathrm{Err}\text{-}5\) corresponding to \(T=1,2,5\) seconds, respectively.

\subsubsection{Orientation Accuracy and Direction-Gated High-Risk Detection}
Let \(\Omega\) denote the set of valid pixels with cardinality \(|\Omega|=N\). For each \(p\in\Omega\), denote the predicted and ground-truth orientations by \(\hat{\theta}(p),\,\theta_{\mathrm{gt}}(p)\in[0,\pi]\), and the predicted and ground-truth TTC by \(\widehat{\mathrm{TTC}}(p)\) and \(\mathrm{TTC}_{\mathrm{gt}}(p)\). We fix an angular tolerance \(\tau_\theta=\pi/12\)  and a TTC threshold \(\tau_T=2\,s\).

% \paragraph{Mean Angular Error (MAE).}
The mean absolute error of orientation aggregates the absolute discrepancy over \(\Omega\):
\begin{equation}
\label{eq:mae_ori}
\mathrm{MAE}_{\mathrm{\theta}}
=\frac{1}{N}\sum_{p\in\Omega}\big|\hat{\theta}(p)-\theta_{\mathrm{gt}}(p)\big|
\end{equation}
For an angular tolerance \(\tau_\theta\), the accuracy within tolerance is
\begin{equation}
\label{eq:acc_theta}
\mathrm{Acc@}\tau_\theta
=\frac{1}{N}\sum_{p\in\Omega}\mathbbm{1}\!\left(\big|\hat{\theta}(p)-\theta_{\mathrm{gt}}(p)\big|\le \tau_\theta\right)
\end{equation}
We define direction-gated high-risk labels by combining a short-TTC condition with an approaching-sector gate:
\begin{equation}
\label{eq:risk_indicators}
\begin{aligned}
y^{\mathrm{gt}}(p) &= \mathbbm{1}\!\big(\mathrm{TTC}_{\mathrm{gt}}(p) < \tau_T \,\land\, \theta_{\mathrm{gt}}(p) \le \tau_\theta\big)\\
\hat{y}(p)         &= \mathbbm{1}\!\big(\widehat{\mathrm{TTC}}(p) < \tau_T \,\land\, \hat{\theta}(p) \le \tau_\theta\big)
\end{aligned}
\end{equation}
Let \(\mathrm{TP},\mathrm{FP},\mathrm{FN}\) be induced by \(\hat{y}\) vs.\ \(y^{\mathrm{gt}}\):
\begin{equation}
\label{eq:confusion}
\begin{aligned}
\mathrm{TP} &= \sum_{p\in\Omega}\mathbbm{1}\!\big(\hat{y}(p)=1 \,\land\, y^{\mathrm{gt}}(p)=1\big)\\
\mathrm{FP} &= \sum_{p\in\Omega}\mathbbm{1}\!\big(\hat{y}(p)=1 \,\land\, y^{\mathrm{gt}}(p)=0\big)\\
\mathrm{FN} &= \sum_{p\in\Omega}\mathbbm{1}\!\big(\hat{y}(p)=0 \,\land\, y^{\mathrm{gt}}(p)=1\big)
\end{aligned}
\end{equation}

We report high-risk precision, high-risk recall, and the high-risk \(F_{\beta}\) score:
\begin{equation}
\label{eq:pr_hrf2}
\begin{aligned}
\mathrm{HR-Precision}   &= \frac{\mathrm{TP}}{\mathrm{TP}+\mathrm{FP}}\\
\mathrm{HR-Recall}    &= \frac{\mathrm{TP}}{\mathrm{TP}+\mathrm{FN}}\\
\mathrm{HR}\text{-}F_{\beta} &= \frac{(1+\beta^2)\,\mathrm{Prec}\cdot \mathrm{Rec}}{\beta^2\,\mathrm{Prec}+\mathrm{Rec}}
\end{aligned}
\end{equation}
we choose \(\beta=2\) to prioritize missed-detection reduction(\(F_2\)).   

\begin{table}[t]
\centering
\caption{Evaluation of MiD and Percentage Errors on nuScenes-SceneFlow}
\label{table:mid_err}
\setlength{\tabcolsep}{4pt} 
\footnotesize               
\begin{adjustbox}{width=\columnwidth,center}
\begin{tabular}{@{}l l c c c c c@{}}
\toprule
\textbf{Input} & \textbf{Method} & \textbf{MiD Err} (\(\downarrow\)) & \textbf{Err-1} (\(\downarrow\)) & \textbf{Err-2} (\(\downarrow\)) & \textbf{Err-5} (\(\downarrow\)) & \textbf{latency(ms)} \\
\midrule
\multirow{4}{*}{Mono}
    & Expansion~\cite{yang2020opticalexpansion} & 79.26           & 0.08          & 15.22         & 18.95         & 67.53 \\
    & BiTTC~\cite{badki2021binaryttc}           & 75.37           & 0.03          & 10.16         & 16.45         & 76.34 \\
    & FP-TTC~\cite{li2024fpttc}                  & 58.68           & 0.04          & 7.84          & 12.69        & \textbf{35.40} \\
    & \textbf{Ours}                             & \textbf{43.23}  & \textbf{0.01} & \textbf{0.47} & \textbf{0.77} & -- \\
\midrule
\multirow{4}{*}{Pano}
    & Expansion~\cite{yang2020opticalexpansion} & 98.04          & 0.33          & 7.02          & 10.19         & 392.40 \\
    & BiTTC~\cite{badki2021binaryttc}           & 86.58          & \textbf{0.10} & 5.43          & 8.82          & 461.60 \\
    & FP-TTC~\cite{li2024fpttc}                  & 70.03          & 0.11          & 5.07          & 8.47          & 255.66 \\
    & \textbf{Ours}                             & \textbf{57.42} & 0.25          & \textbf{2.04} & \textbf{3.33} & \textbf{102.32} \\
\bottomrule
\end{tabular}
\end{adjustbox}
\end{table}

\begin{table}[t]
\centering
\caption{Frame-level Recall and False Alarm Rate (FAR) on Real Vehicle Platform}
\label{table:local_recall_far}
\setlength{\tabcolsep}{3pt} 
\small 
\begin{adjustbox}{max width=\columnwidth,center}
\begin{tabular}{@{}l l c c@{}}
\toprule
\textbf{Input} & \textbf{Method} & \textbf{Recall} (\(\uparrow\)) & \textbf{FAR} (\(\downarrow\)) \\
\midrule
\multirow{4}{*}{Pano}
&Expansion~\cite{yang2020opticalexpansion} & 0.863          & 0.063   \\
&BiTTC~\cite{badki2021binaryttc}           & 0.917          & 0.031   \\
&FP-TTC~\cite{li2024fpttc}                 & 0.955          & 0.020   \\
&\textbf{Ours}                             & \textbf{0.976} & \textbf{0.011} \\
\bottomrule
\end{tabular}
\end{adjustbox}

\vspace{2pt}
{\footnotesize
\textit{Note:} Recall is computed over frames with ground-truth collision risk,
while FAR is computed over frames without ground-truth collision risk.
}
\end{table}


\begin{figure}[!t]  
    \centering
    \includegraphics[width=0.48\textwidth]{figs/hardware.pdf}
    \caption{(a), (b) Real vehicle platform and surround-view cameras. (c) Camera orientations and overlapping fields of view.}
    \label{fig:hareware}
\end{figure}

\subsection{Comparison with Monocular Baselines}

Since existing TTC methods are monocular, we design single-view and multi-view comparison protocols. In the single-view setting, we use the camera-view masks generated during ground-truth construction to extract the front-view region from our panoramic predictions and compare it directly with monocular baselines on the same front camera. In the multi-view setting, we apply each monocular baseline independently to all surround-view cameras and then compare their aggregated performance against our panoramic model. 

As shown in Table~\ref{table:mid_err}, our method consistently outperforms all monocular baselines in the single-view setting. 
Compared with the strongest baseline FP-TTC, SCOPE reduces MiD Err from 58.68 to 43.23 and substantially lowers the percentage errors at loose thresholds Err-T. This demonstrates that our panoramic training and OTTC formulation improve both the accuracy and the stability of TTC estimates, even when they are evaluated only on the front camera.

The advantage becomes more pronounced in the multi-view setting. 
When monocular methods are applied independently to all surround-view cameras, their MiD Err increases, such as FP-TTC rising from 58.68 to 70.03. 
This behavior reflects inconsistencies across views and duplicated or missing predictions in overlapping regions. 
In contrast, our unified panoramic model achieves the lowest MiD Err of 57.42 and lower Err-2 and Err-5, while requiring only 102.32\,ms per surround-view frame. 
This latency is about 2.5 times lower than that of the strongest monocular baseline when applied independently to all surround-view cameras.
These results show that a single RV-based panoramic predictor is both more accurate and more efficient than aggregating multiple monocular TTC estimators.

\begin{figure*}[!t]  
    \centering
    \includegraphics[width=0.98\textwidth]{figs/cyberrock_result0.pdf}
    \caption{Visualization of zero-shot performance on our real vehicle platform. The figure displays front (top row) and back (bottom row) collision scenes. Columns from left to right show: the surround-view RGB inputs, the estimated OTTC (TTC Value and Orientation), and the detected collision points projected onto the surround-view RGB images. The results demonstrate the model's capability to generalize to real-world scenarios without specific training.}
    \label{fig:campus_zeroshot}
\end{figure*}



\subsection{Zero-Shot Generalization to Campus Scenes}
\label{sec:zeroshot_campus}
Our model is trained exclusively on the proposed \textit{nuScenes-SceneFlow} dataset. Using this trained model, we can deploy it on our real vehicle platform in a purely zero-shot manner. By replacing the camera intrinsics and extrinsics with those of the local sensor setup, we obtain panoramic TTC value and orientation maps without any fine-tuning. 

To quantitatively evaluate the generalization performance of the model, we tested different models on real vehicle platform. We evaluate collision warning performance at the frame level using two complementary metrics: \textbf{Recall} and \textbf{False Alarm Rate (FAR)}. Recall measures how often the system successfully issues a warning on frames that truly contain collision risk. FAR measures how often the system raises a warning on frames that are in fact collision-free.

Following the TTC4MCP~\cite{li2023ttc4mcp} framework, we apply joint thresholds on the predicted TTC value and orientation maps to extract collision risk regions, and then represent each region by a collision point for visualization and evaluation. Frames that satisfy the collision condition are treated as positive frames. A positive frame is counted as a true positive (TP) if at least one predicted collision point falls within a spatial tolerance region around the annotated collision area; otherwise, it is counted as a false negative (FN). The frame-level recall is computed as TP\(_\text{frame}\) / (TP\(_\text{frame}\) + FN\(_\text{frame}\)).
Frames without collision risk are treated as negative frames. If any collision point is predicted in such a frame, it is counted as a false alarm; otherwise, it is a true negative. The frame-level FAR is defined as the ratio of false alarms to all negative frames. 
% Quantitative results on our real vehicle platform are summarized in Table~\ref{table:local_recall_far}, while qualitative comparisons between FP-TTC~\cite{li2024fpttc} and our method are illustrated in Fig.~\ref{fig:campus_zeroshot}.

Fig.~\ref{fig:campus_zeroshot} and Table~\ref{table:local_recall_far} collectively demonstrate the impressive zero-shot performance of our method in real-world collision scenarios. Compared with the strongest baseline, FP-TTC, SCOPE achieves a notable improvement in safety reliability, increasing frame-level recall from 0.955 to 0.976 and decreasing FAR from 0.020 to 0.011. While other baselines like Expansion and BiTTC suffer from either missed detections or excessive false alarms, our panoramic OTTC model achieves a superior trade-off, maintaining high sensitivity to collision risks while effectively suppressing false positives. These results confirm that our proposed range view representation and orientation-aware risk modeling generalize robustly from nuScenes to unseen real-world environments without requiring retraining.

% The qualitative results in Fig.~\ref{fig:campus_zeroshot} demonstrate that our method achieves impressive zero-shot performance on real collision scenarios. 
% From Table~\ref{table:local_recall_far}, we observe that our method achieves the best performance on our real vehicle platform. 
% Compared with the strongest baseline FP-TTC, SCOPE improves frame-level recall from 0.955 to 0.976 and reduces FAR from 0.020 to 0.011. 
% Expansion and BiTTC either miss more risky frames or trigger more false alarms, while our panoramic OTTC model maintains high sensitivity to true collision cases and simultaneously suppresses false alarms. 
% These results indicate that the proposed range view training and orientation-aware risk representation generalize well from nuScenes to the unseen environment in a purely zero-shot setting.

% The qualitative results in Fig.~\ref{fig:campus_zeroshot} further explain these gains. 
% When FP-TTC is applied independently to each surround-view image, it produces inconsistent TTC fields in overlapping regions and often highlights static background as risky, which leads to redundant or incorrect collision area. 
% In contrast, our method fuses information from all cameras in the unified range-view domain and yields a globally consistent TTC field. 
% Collision points selected from our TTC value and orientation maps concentrate on truly approaching agents, such as pedestrians and vehicles that move toward the ego car, while non-threatening objects are largely suppressed. 
% This behavior demonstrates that SCOPE not only transfers to new scenes without fine-tuning, but also provides cleaner and more robust collision warnings in real driving scenarios.




\begin{table*}[t]
\centering
\caption{Ablation: Depth Sampling Strategy on nuScenes. Acc@\(\tau_\theta\) uses \(\tau_\theta=\pi/12\) rad; HR-\(F_2\) uses \(\beta=2\).}
\label{tab:ablation_depth}
\setlength{\tabcolsep}{4pt}
\footnotesize
\begin{adjustbox}{width=\textwidth,center}
\begin{tabular}{@{}l cccc ccccc@{}}
\toprule
& \multicolumn{4}{c}{\textbf{Scale metrics}} & \multicolumn{5}{c}{\textbf{Orientation metrics}} \\
\cmidrule(lr){2-5} \cmidrule(lr){6-10}
\textbf{Variant} 
& \textbf{MiD Err} (\(\downarrow\)) & \textbf{Err-1} (\(\downarrow\)) & \textbf{Err-2} (\(\downarrow\)) & \textbf{Err-5} (\(\downarrow\))
& \textbf{MAE\(_\theta\)}(\(\downarrow\)) & \textbf{Acc@\(\tau_\theta\)} (\(\uparrow\)) & \textbf{Precision}(\(\uparrow\)) & \textbf{Recall}(\(\uparrow\)) & \textbf{HR-\(F_2\)} (\(\uparrow\)) \\
\midrule
Linear, \(K=4\)        & 82.29 & \textbf{0.22} & 2.96 & 4.90 & 0.07 & 0.98 & 0.71 & \textbf{0.91} & 0.86 \\
Linear, \(K=8\)        & 69.83 & 0.28 & 2.24 & 3.94 & 0.05 & \textbf{0.99} & 0.84 & 0.85 & 0.85 \\
Linear, \(K=16\)       & 82.78 & 0.23 & 3.42 & 4.31 & 0.06 & 0.98 & 0.82 & 0.86 & 0.85 \\
Log-uniform, \(K=4\)   & 66.96 & 0.23 & 2.59 & 3.99 & 0.05 & \textbf{0.99} & \textbf{0.88} & 0.87 & 0.87 \\
\textbf{Log-uniform, \(\mathbf{K{=}8}\)}
    & \textbf{57.42} & 0.25 & \textbf{2.04} & \textbf{3.33} & \textbf{0.04} & \textbf{0.99} & \textbf{0.88} & 0.88 & \textbf{0.88} \\
Log-uniform, \(K=16\)  & 61.57 & 0.25 & 2.27 & 3.68 & \textbf{0.04} & 0.98 & 0.87 & 0.87 & 0.87 \\
\bottomrule
\end{tabular}
\end{adjustbox}
\end{table*}

\begin{table*}[t]
\centering
\caption{Ablation: Source of Query Initialization on nuScenes. Acc@\(\tau_\theta\) uses \(\tau_\theta=\pi/12\) rad; HR-\(F_2\) uses \(\beta=2\); risk gate uses \(\tau_T=2\) s.}
\label{tab:ablation_query}
\setlength{\tabcolsep}{4pt}
\footnotesize
\begin{adjustbox}{width=\textwidth,center}
\begin{tabular}{@{}l cccc ccccc@{}}
\toprule
& \multicolumn{4}{c}{\textbf{Scale metrics}} & \multicolumn{5}{c}{\textbf{Orientation metrics}} \\
\cmidrule(lr){2-5} \cmidrule(lr){6-10}
\textbf{Variant}
& \textbf{MiD Err} (\(\downarrow\))
& \textbf{Err-1} (\(\downarrow\))
& \textbf{Err-2} (\(\downarrow\))
& \textbf{Err-5} (\(\downarrow\))
& \textbf{MAE\(_\theta\)} (\(\downarrow\))
& \textbf{Acc@\(\tau_\theta\)} (\(\uparrow\))
& \textbf{Precision} (\(\uparrow\))
& \textbf{Recall} (\(\uparrow\))
& \textbf{HR-\(F_2\)} (\(\uparrow\)) \\
\midrule
Zero-Init            & 104.39 & 1.72 & 5.14 & 5.68 & 0.10 & 0.92 & 0.63 & 0.80 & 0.75 \\
ConcatFeat-Init     & 71.78  & 0.26 & 2.64 & 3.95 & 0.05 & 0.98 & 0.84 & 0.85 & 0.85 \\
\textbf{GVB-Init}  & \textbf{57.42}  & \textbf{0.25} & \textbf{2.04} & \textbf{3.33} &
\textbf{0.04} & \textbf{0.99} & \textbf{0.88} & \textbf{0.88} & \textbf{0.88} \\
\bottomrule
\end{tabular}
\end{adjustbox}
\end{table*}

\begin{table*}[t]
\centering
\caption{Ablation: Sampling Mechanism on nuScenes. Acc@\(\tau_\theta\) uses \(\tau_\theta=\pi/12\) rad; HR-\(F_2\) uses \(\beta=2\); risk gate uses \(\tau_T=2\) s.}
\label{tab:ablation_mechanism}
\setlength{\tabcolsep}{4pt}
\footnotesize
\begin{adjustbox}{width=\textwidth,center}
\begin{tabular}{@{}l cccc ccccc@{}}
\toprule
& \multicolumn{4}{c}{\textbf{Scale metrics}} & \multicolumn{5}{c}{\textbf{Orientation metrics}} \\
\cmidrule(lr){2-5} \cmidrule(lr){6-10}
\textbf{Variant}
& \textbf{MiD Err} (\(\downarrow\))
& \textbf{Err-1} (\(\downarrow\))
& \textbf{Err-2} (\(\downarrow\))
& \textbf{Err-5} (\(\downarrow\))
& \textbf{MAE\(_\theta\)} (\(\downarrow\))
& \textbf{Acc@\(\tau_\theta\)} (\(\uparrow\))
& \textbf{Precision} (\(\uparrow\))
& \textbf{Recall} (\(\uparrow\))
& \textbf{HR-\(F_2\)} (\(\uparrow\)) \\
\midrule
Geometry-only
& 63.57 & 0.28 & 2.61 & 4.05 & 0.06 & 0.95 & 0.83 & 0.84 & 0.84 \\
Learning-only
& 59.17 & \textbf{0.25} & 2.13 & \textbf{3.28} & 0.05 & \textbf{0.99} & 0.86 & 0.86 & 0.86 \\
\textbf{Geometry + Refinement}
& \textbf{57.42} & \textbf{0.25} & \textbf{2.04} & 3.33 & \textbf{0.04} & \textbf{0.99} & \textbf{0.88} & \textbf{0.88} & \textbf{0.88} \\
\bottomrule
\end{tabular}
\end{adjustbox}
\end{table*}

\begin{table*}[t]
\centering
\caption{Ablation: Camera-Level Embedding on nuScenes. Acc@\(\tau_\theta\) uses \(\tau_\theta=\pi/12\) rad; HR-\(F_2\) uses \(\beta=2\); risk gate uses \(\tau_T=2\) s.}
\label{tab:ablation_camembed}
\setlength{\tabcolsep}{4pt}
\footnotesize
\begin{adjustbox}{width=\textwidth,center}
\begin{tabular}{@{}l cccc ccccc@{}}
\toprule
& \multicolumn{4}{c}{\textbf{Scale metrics}} & \multicolumn{5}{c}{\textbf{Orientation metrics}} \\
\cmidrule(lr){2-5} \cmidrule(lr){6-10}
\textbf{Variant}
& \textbf{MiD Err} (\(\downarrow\))
& \textbf{Err-1} (\(\downarrow\))
& \textbf{Err-2} (\(\downarrow\))
& \textbf{Err-5} (\(\downarrow\))
& \textbf{MAE\(_\theta\)} (\(\downarrow\))
& \textbf{Acc@\(\tau_\theta\)} (\(\uparrow\))
& \textbf{Precision} (\(\uparrow\))
& \textbf{Recall} (\(\uparrow\))
& \textbf{HR-\(F_2\)} (\(\uparrow\)) \\
\midrule
w/o (no camera ID)
& 60.85 & 0.27 & 2.20 & 3.37 & 0.05 & 0.98 & 0.87 & 0.86 & 0.87 \\
\textbf{w/ (learnable ID)}
& \textbf{57.42} & \textbf{0.25} & \textbf{2.04} & \textbf{3.33} & \textbf{0.04} & \textbf{0.99} & \textbf{0.88} & \textbf{0.88} & \textbf{0.88} \\
\bottomrule
\end{tabular}
\end{adjustbox}
\end{table*}

 

\subsection{Ablation Study}
\label{sec:ablation}


\textbf{Depth Sampling Strategy.}
We keep the depth range \([1,40]\)m and the RVT configuration unchanged, while modifying the depth distribution of sampling points and the number of sampling points. Concretely, we evaluate:
\begin{enumerate}
  \item Linear-uniform sampling with \(K\in\{4,8,16\}\): bins are evenly spaced in metric depth.
  \item Log-uniform sampling with \(K\in\{4,8,16\}\): bins are evenly spaced in log-depth, allocating more points to the far field.
\end{enumerate}

The results show that varying the number of depth samples \(K\) and their distribution significantly impacts performance. When keeping the distribution fixed, using \(K=4\) often misses fine details due to under-sampling, while \(K=8\) achieves a good balance between accuracy and computational cost. Increasing to \(K=16\) provides limited improvement and may even introduce noise or redundancy, increasing computational overhead. Comparing linear and log-uniform distributions at a fixed \(K\), log-uniform sampling consistently improves far-range accuracy by allocating more samples to distant regions where depth changes more gradually, while linear sampling slightly benefits very close-range accuracy but performs worse at greater distances. Overall, the best trade-off is achieved with log-uniform sampling and \(K=8\), which balances accuracy and efficiency effectively. Smaller \(K\) is faster but risks losing depth details, while larger \(K\) increases computational demands without significant benefits. The results are shown in Table~\ref{tab:ablation_depth}.


\textbf{Query Initialization Strategy.}
We use DepthAnything (DA) as a teacher to guide the RVT module. The teacher's guidance is gradually reduced during training to remove external depth reliance at inference. We compare the following methods:
\begin{enumerate}
    \item Zero initialization (Zero-Init): The student query is initialized to zero after light convolution and positional encoding. This tests the ability to recover from a neutral prior.
    \item Concatenate feature initialization (ConcatFeat-Init): The student query is initialized by concatenating per-camera features along the width. The per-camera features are easy to obtain and contain rich appearance information, but may introduce seams at image boundaries due to lack of geometric consistency.
    \item Use DA-supervised geometric visibility bootstrap initialization (GVB-Init): we initialize the student query by ray casting with visibility masks and a single deformable aggregation, resulting in a smooth and seam-free prior under DA supervision.
\end{enumerate}

The results are shown in Table~\ref{tab:ablation_query}. The evaluation results demonstrate that the Learnable Feat Initialization approach achieves the best overall performance across all metrics. Compared to ConcatFeat-Init, it significantly reduces MiD Err from 71.78 to 57.42, lowers Err-2 and Err-5 from 2.64 to 2.04 and 3.95 to 3.33, respectively, and improves HR-\(F_2\) from 0.85 to 0.88, while maintaining the highest Acc@\(\tau_\theta\) of 0.99. In contrast, Zero-Init performs poorly, with MiD Err increasing to 104.39, orientation accuracy dropping to an Acc@\(\tau_\theta\) of 0.92, and HR-\(F_2\) falling to 0.75, indicating that a neutral initialization is difficult to optimize effectively, even with early guidance from DA.

These findings suggest that while DA guidance helps stabilize early training, the choice of initialization strategy plays a critical role in the final performance. ConcatFeat-Init provides a significant gain over the Zero-Init baseline. Nevertheless, its dependence on per-camera features induces boundary artifacts, resulting in a slight drop in recall and HR-\(F_2\). In contrast, GVB-Init provides a smooth, geometry-aligned prior that eliminates seams and consistently enhances both TTC value and orientation metrics. For practical applications, GVB-Init is recommended for its superior accuracy and stability, while ConcatFeat-Init serves as a viable mid-tier alternative for simpler implementations. Zero-Init, though functional, is not ideal for deployment due to its lower performance. 

% \begin{figure*}[t]
%     \centering
%     \begin{minipage}{0.98\textwidth}
%         \centering
%         \includegraphics[width=\linewidth]{figs/cyberrock_result0.pdf}
%     \end{minipage}

%     \smallskip 

%     \begin{minipage}{0.98\textwidth}
%         \centering
%         \includegraphics[width=\linewidth]{figs/cyberrock_result1.pdf}
%     \end{minipage}
%     \caption{Zero-shot performance on collision scenarios from our local autonomous driving platform. 
%     The first scenario corresponds to a forward collision, while the second scenario corresponds to a rear collision.
%     From top to bottom are 
%     (a) surround-view RGB images,
%     (b) monocular scale predictions of FP-TTC~\cite{li2024fpttc} applied independently to each camera, and
%     (c) range view scale predictions of our panoramic model. 
%     \ding{72} denotes collision points obtained from our OTTC field by jointly thresholding scale and orientation. 
%     \protect\legendbox{lightblue}/\protect\legendbox{darkred}: monocular FP-TTC yields inconsistent risk estimates in overlapping regions, assigning different TTC values to the same physical object across views, whereas our method fuses all cameras in the range-view domain and produces a single, globally consistent panoramic TTC field. 
%     \protect\legendbox{lightgreen}: FP-TTC marks a safe pedestrian and a parked vehicle as risky incorrectly, while our orientation-aware method isolates the only real threat, resulting in cleaner and more precise collision warnings.}
%     \label{fig:campus_zeroshot}
% \end{figure*}



% \textbf{Sampling Mechanism: Geometry vs.\ Learning.}
% We study how sampling locations and aggregation weights are chosen inside the RVT, under three configurations:
% \begin{enumerate}
%   \item \emph{Geometry-only}: image-plane samples are computed from camera intrinsics/extrinsics using ray casting over depth bins with explicit visibility masking. Weights are proportional to visibility (masked and renormalized). No learnable offsets or weights are used.
%   \item \emph{Learning-only}: both sampling locations and weights are predicted from the query content by the network (deformable attention heads), without geometric anchors or visibility masks. This variant relies purely on end-to-end learning.
%   \item \emph{Geometry + refinement}: geometric samples serve as anchors; the network predicts small residual offsets on top of them, and refines weights under visibility masking with per-head softmax normalization. This combines hard geometry with data-driven flexibility.
% \end{enumerate}

% The results are shown in Table~\ref{tab:ablation_mechanism}. Against the geometry baseline, Geometry+Refinement lowers the MiD error by roughly 7\% and reduces Err-5 by about 19\%, indicating that geometric anchors plus small learned offsets help most on hard regions such as boundaries and occlusions. Learning-only pushes the MiD error down further—about 10\% vs.\ geometry and 3\% vs.\ refinement—and yields the strongest orientation quality (e.g., $\sim\!\mathbf{33\%}$ lower angular MAE vs.\ geometry and $\sim\!\mathbf{20\%}$ vs.\ refinement). It also improves detection-style metrics: precision increases by about 4--6\% and recall by 2--5\% relative to geometry, with a comparable boost in HR-$F_2$. In practice, geometry+refinement is a robust default when stability near view edges and occlusions is critical; learning-only is preferred when the goal is the strongest overall accuracy.

\textbf{Sampling Mechanism: Geometry vs.\ Learning.}
We study how sampling locations and aggregation weights are chosen inside the RVT, under three configurations:
\begin{enumerate}
  \item \emph{Geometry-only}: image-plane samples are computed from camera intrinsics/extrinsics using ray casting over depth bins with explicit visibility masking. Weights are proportional to visibility (masked and renormalized). No learnable offsets or weights are used.
  \item \emph{Learning-only}: both sampling locations and weights are predicted from the query content by the network (deformable attention heads), without geometric anchors or visibility masks. This variant relies purely on end-to-end learning.
  \item \emph{Geometry + refinement}: geometric samples serve as anchors; the network predicts small residual offsets on top of them, and refines weights under visibility masking with per-head softmax normalization. This combines hard geometry with data-driven flexibility.
\end{enumerate}

% As shown in Table~\ref{tab:ablation_mechanism}, both learned variants improve over Geometry-only, and \emph{Geometry + refinement} achieves the best overall performance. It reduces MiD error from 63.57 (Geometry-only) to 57.42, and further improves orientation accuracy, lowering MAE\(_\theta\) from 0.06 to 0.04 and raising HR-\(F_2\) from 0.84 to 0.88. Compared with Learning-only, Geometry + refinement provides consistently better scale and orientation metrics while retaining high Acc@\(\tau_\theta\), indicating that geometric anchors plus modest learned corrections offer the most effective trade-off between stability and accuracy.

As detailed in Table~\ref{tab:ablation_mechanism}, both learned variants outperform the Geometry-only baseline, with \emph{Geometry + refinement} emerging as the most robust method. It not only reduces the baseline MiD error by a significant margin (63.57 \(\to\) 57.42) and improves orientation metrics (MAE\(_\theta\): 0.06 \(\to\) 0.04; HR-\(F_2\): 0.84 \(\to\) 0.88), but also demonstrates superior consistency over the Learning-only model. This confirms that leveraging geometric anchors alongside learned corrections establishes the most effective trade-off between stability and accuracy.


\textbf{Camera-Level Embedding.}
We test whether an explicit camera identity helps cross-view fusion and stability. We consider two configurations:
\begin{enumerate}
    \item \emph{Without Camera ID}: No camera-specific information is added to the features; the RVT processes only the appearance features from each view.
    \item \emph{With Learnable Camera ID}: A unique, trainable embedding is added to the features of each camera before fusion. This helps the model handle differences in field of view, mounting position, and color/exposure settings for each camera.
\end{enumerate}
We try to figure out whether identity priors improve cross-view alignment and reduce boundary inconsistencies.

% The results are shown in Table~\ref{tab:ablation_camembed}. Relative to the no-ID baseline, MiD error drops by roughly 5--6\%, small-error rates improve by 7\%, and the long-tail error shrinks slightly. Orientation improves more notably: angular MAE decreases by about 20\%, and detection-style metrics (accuracy, precision, recall, HR-$F_2$) increase by around 1--2 percentage points. The benefits are most pronounced for orientation and recall, suggesting that camera identity helps the model resolve ambiguous matches near view boundaries and under photometric shifts. A learnable camera ID is a low-cost prior that consistently improves fusion quality and training stability; we therefore enable it by default.

% The results are summarized in Table~\ref{tab:ablation_camembed}. Relative to the no-ID baseline, MiD error drops by roughly 5--6\% and the short-range metrics \text{Err-1,2} improves by about 7\%. Orientation benefits even more. Angular MAE decreases by around 20\%, while \text{Acc@}(\tau_\theta), \text{Prec@}(\tau_\theta), \text{Recall@}(\tau_\theta), and HR-\(F2\) all increase by approximately 1--2 percentage points. These gains suggest that camera identity helps the model resolve ambiguous matches near view boundaries and under photometric shifts. As a low-cost prior that improves fusion quality and training stability, the learnable camera ID is therefore enabled by default.

The results are summarized in Table~\ref{tab:ablation_camembed}. Relative to the no-ID baseline, MiD error drops by roughly 5--6\% and the short-range metrics $\text{Err-1}$ and $\text{Err-2}$ improve by about 7\%. Orientation benefits even more. Angular MAE decreases by around 20\%, while $\text{Acc@}(\tau_\theta)$, $\text{Prec@}(\tau_\theta)$, $\text{Recall@}(\tau_\theta)$, and HR-$F_2$ all increase by approximately 1--2 percentage points. These gains suggest that camera identity helps the model resolve ambiguous matches near view boundaries and under photometric shifts. As a low-cost prior that improves fusion quality and training stability, the learnable camera ID is therefore enabled by default.


\section{Conclusion}
\label{sec:con}

We presented \textbf{SCOPE}, a surround-view collision prediction framework that estimates \textbf{OTTC} on a unified spherical range view domain, enabling omnidirectional and direction-aware collision assessment from surround-view inputs. 
SCOPE projects multi-view features into a continuous RV domain and jointly regresses TTC with relative motion orientation. Through this design, it addresses the \textbf{spatial incompleteness} of forward-only baselines and the \textbf{directional ambiguity} of scalar TTC metrics. The final output is a geometrically consistent panoramic OTTC field that downstream modules can consume directly.
We further constructed a densified \textit{nuScenes-SceneFlow} dataset, which generates pixel-wise TTC value and orientation supervision in the range view domain. 
Thorough experiments on public benchmarks demonstrate that SCOPE achieves SOTA MiD and TTC accuracy and reducing false alarms. Furthermore, the system supports zero-shot deployment on a real vehicle platform while satisfying real-time latency constraints.

Comprehensive ablations clarified key design choices. For depth sampling, log-uniform bins with a moderate number of samples offer a good trade-off between depth coverage and efficiency. For sampling mechanisms inside the Range View Transformer, combining geometric anchors with small learned refinements achieves the best overall performance. A lightweight camera embedding enhances fusion stability and improves the performance of OTTC estimation.

In future work, we plan to incorporate longer temporal contexts to enhance prediction stability. Additionally, we will explore self-supervised and semi-supervised learning schemes to mitigate the reliance on dense LiDAR-derived labels, and finally, evaluate the framework within closed-loop planning and control stacks.
We conclude that panoramic OTTC estimation provides a solid foundation for the safe and scalable deployment of autonomous mobile robots in complex, dynamic environments.


%%%%%%%%%%%%%%%%%%%%%%%%%%%%%%%%%%%%%%%%%%%%%%%%%%%%%%%%%%%%%%%%%%%%%%%%%%%%%%%%
%%%%%%%%% REFERENCES
{\small
\bibliographystyle{IEEEtran}
\bibliography{refs}
\vspace{-50mm}
}

\newpage

% \begin{IEEEbiography}
% [{\includegraphics[width=1in,height=1.25in,clip,keepaspectratio]{figs/bio/NanMing.jpg}}]{Nan Ming}
% received the B.S. degree in artificial intelligence from Tongji University, Shanghai, China, in 2023. He is currently pursuing his M.S. degree with the
% Department of Automation at Shanghai Jiao Tong University.\\
% His main research interests include computer vision, machine learning, 3D scene understanding, and their applications in intelligent transportation systems.
% \vspace{-50 mm}
% \end{IEEEbiography}

% \begin{IEEEbiography}
% [{\includegraphics[width=1in,height=1.25in,clip,keepaspectratio]{figs/bio/YeqiangQian.jpg}}]{Yeqiang Qian}
% received a Ph.D. degree in control science and engineering from Shanghai Jiao Tong University, Shanghai, China, in 2020. He is currently a Tenure Track Associate Professor with the Department of Automation at Shanghai Jiao Tong University, Shanghai, China. \\
% His main research interests include computer vision, pattern recognition, machine learning, and their applications in intelligent transportation systems.
% \vspace{-50 mm}
% \end{IEEEbiography}

% \begin{IEEEbiography}
% [{\includegraphics[width=1in,height=1.25in,clip,keepaspectratio]{figs/bio/ChunYuFeng.jpg}}]{Chunyu Feng}
% received his B.S. degree in Automation in 2019 and his M.S. degree in Automation in 2022, both from Tongji University, Shanghai, China. He is currently working toward the Ph.D. degree in Control Science and Engineering at Shanghai Jiao Tong University.\\
% His main research interests include computer vision, AEB system design, and their applications in intelligent transportation systems. 

% \vspace{-50 mm}
% \end{IEEEbiography}

% \begin{IEEEbiography}[{\includegraphics[width=1in,height=1.25in,clip,keepaspectratio]
% {figs/bio/ChunxiangWang.jpg}}]{Chunxiang Wang}
% received the Ph.D. degree in mechanical engineering from the Harbin Institute of Technology, Harbin, China, in 1999.\\
% She is currently an Associate Professor with the Department of Automation at Shanghai Jiao Tong University, Shanghai, China. Her research interests include robotic technology and electromechanical integration.
% \vspace{-50 mm}
% \end{IEEEbiography}

% \begin{IEEEbiography}[{\includegraphics[width=1in,height=1.25in,clip,keepaspectratio]
% {figs/bio/MingYang.jpg}}]{Ming Yang}
% received the Master and Ph.D. degrees from Tsinghua University, Beijing, China, in 1999 and 2003, respectively.\\
% He is currently the Full Tenure Professor at Shanghai Jiao Tong University, the deputy director of the Innovation Center of Intelligent Connected Vehicles. He has been working in the field of intelligent vehicles for more than 20 years. He participated in several related research projects, such as the THMRV project (first intelligent vehicle in China), European CyberCars and CyberMove projects, CyberC3 project, CyberCars-2 project, ITER transfer cask project, AGV, etc.
% \vspace{-50 mm}
% \end{IEEEbiography}


\end{document}


